%% toc: P15 171 171.11
\chapter{Generalization Failure under High Capacity or Distribution Shift}
\label{sec:p15-w02b-p15-066}

A sampling guarantee for a learned predictor rests on two separate
hypotheses. The first concerns the sample: the selection rule chooses
from a class fixed in advance whose capacity is small relative to the
number of observations. The second concerns the future: the law under
which the predictor will be used is the law that produced the sample.
\Cref{sec:p15-bounded-concentration-learning,sec:p15-occam-learning,sec:p15-vc-dimension-and-uniform-learning}
prove guarantees under both hypotheses and record several isolated
cases in which one of them fails. This chapter treats the two failures
as a subject in their own right. For high capacity it identifies, for
each earlier capacity counterexample, the fixed class that actually
contains every output the selection rule can return, and shows that the
counterexample is an instance of a class too rich for the sample used. For distribution
shift it proves how a guarantee obtained under the training law
transfers to a changed law, shows by examples that each transfer
inequality is attained, and proves that no amount of training data
from the old law can replace an assumption linking the two laws.

\emph{Inputs.} A measurable space of records, a training law \(P\),
an evaluation law \(Q\) on the same space, a loss with values in
\([0,1]\), a candidate class, a selection rule and an independent
training sample from \(P\).
\emph{Conclusions.} For each earlier capacity failure, the fixed class that
actually contains every possible output and the reason its guarantee is
uninformative; transfer inequalities from \(P\) to \(Q\) in terms of
total variation and of a bounded density ratio, each with an attaining
example; and the impossibility of a certificate for a changed law that
uses the training sample alone.

The framework sentence that training and test laws coincide, and the
remark that they can differ, are those of von Luxburg and
Sch\"olkopf's list of assumptions
\cite[Section~2.2, pp.~3--4]{LuxburgScholkopf2008Author}, whose
memorizing classifier \cite[Section~4.2, pp.~13--14]{LuxburgScholkopf2008Author}
is already treated in \cref{thm:p15-vcl-examples}. Their no-free-lunch
discussion \cite[Section~8, pp.~30--32]{LuxburgScholkopf2008Author}
states, as Theorem~9 on p.~32, that for each sample size and each
classifier some law of Bayes risk zero gives expected risk above one
half minus any fixed margin; they state it with reference to other
works and do not prove it there. That theorem is not proved in this
book; the weaker lower bounds inside a class of large dimension that the
book does prove are those of \cref{sec:p15-w02b-p15-065-lower}. The
results of this chapter are the book's own and are proved in the annex.

\section{The two hypotheses behind a sampling guarantee}
\label{sec:p15-w02b-p15-066-hypotheses}

\begin{hbdefinition}{def:p15-w02b-p15-066-two-laws}{Two laws and an enclosing class}
Let \((\mathcal Z,\mathcal B)\) be a measurable space of records and let
\(P\) and \(Q\) be probability measures on it: \(P\) is the
\emph{training law}, which produces an independent sample
\(Z_1,\ldots,Z_n\), \(n\geq1\), and \(Q\) is the \emph{evaluation law},
under which a chosen predictor is used. For a candidate \(h\) with
measurable loss \(\ell(h,\cdot):\mathcal Z\to[0,1]\), put
\[
 R_P(h)=\int\ell(h,z)\,P(dz),\qquad
 R_Q(h)=\int\ell(h,z)\,Q(dz),
\]
\[
 \widehat R_n(h)=\frac1n\sum_{i=1}^n\ell(h,Z_i).
\]
Thus \(R_P\) is the risk \(R\) of \cref{def:p15-conc-learning}; the
second argument of \(R_P(h,f)\) in \cref{def:p15-occam-model} is the
target, which here is absorbed into the loss. The distance between the
two laws is the event-supremum total variation
\(\|P-Q\|_{\mathrm{TV}}\) of \cref{def:p15-weak-total-variation}.
\emph{Distribution shift} means \(Q\neq P\).

When records are pairs \(z=(x,y)\) with a finite label set, \(Q\)
differs from \(P\) by \emph{covariate shift} when the two laws have the
same conditional label probabilities \(p(y\mid x)\) and different input
marginals \(P_{\mathcal X},Q_{\mathcal X}\); precisely, there are
measurable \(p(y\mid\cdot)\geq0\) with \(\sum_yp(y\mid x)=1\) and, for
every measurable \(C\subseteq\mathcal X\times\mathcal Y\),
\[
 \begin{aligned}
 P(C)&=\int\sum_{y}\mathbf1_C(x,y)\,p(y\mid x)\,P_{\mathcal X}(dx),\\
 Q(C)&=\int\sum_{y}\mathbf1_C(x,y)\,p(y\mid x)\,Q_{\mathcal X}(dx).
 \end{aligned}
\]
It differs by \emph{label shift} or \emph{concept change} when the
conditional label probabilities themselves change.

A selection rule maps the sample, and possibly independent internal
randomness, to a predictor. Its \emph{enclosing class} is a class fixed
before sampling that contains every predictor the rule can return.
The finite-class, weighted and capacity guarantees of
\cref{thm:p15-conc-finite-learning,thm:p15-occam-finite,thm:p15-occam-weighted,thm:p15-vcl-erm}
apply to the selected predictor only through such a class, and only
for the risk under the training law.
\index{distribution shift}\index{covariate shift}\index{generalization!failure}
\end{hbdefinition}

A guarantee can therefore fail in two distinct ways. It fails by
\emph{capacity} when the smallest enclosing class is so rich that its
uniform deviation is not small at the sample size used; the training
error of the selected predictor then says little about its risk even
under \(P\). It fails by \emph{shift} when \(Q\neq P\): the guarantee
concerns \(R_P\) and is silent about \(R_Q\) until some assumption
links \(Q\) to \(P\). The two failures are independent of each other.
\Cref{thm:p15-occam-boundaries}, item~1, records a third, a failure of
independent sampling, which this chapter does not treat further.

\section{High capacity: fits that certify nothing}
\label{sec:p15-w02b-p15-066-capacity}

The upper bounds of \cref{thm:p15-vcl-samples} show that finite
dimension suffices for learning at every law. The matching necessity is
proved in \cref{sec:p15-w02b-p15-065-lower}: a class that shatters many
points defeats every learner, whatever its selection rule, on some law
under which the class contains a perfect predictor
(\cref{thm:p15-w02b-p15-065-lower-realizable}), and a class of infinite
dimension admits no distribution-free sample size at all
(\cref{thm:p15-w02b-p15-065-lower-infinite}). The lower bound of
\cref{thm:p15-occam-unseen-labels} is the instance of the same averaging
argument for the class of all binary functions on a finite domain.

The four capacity counterexamples proved earlier are now seen to share
one cause: in each, the predictor returned after seeing the sample lies
in no fixed class small enough for its sample size.

\begin{hbtheorem}{thm:p15-w02b-p15-066-enclosing}{proved=annex:p15-w02b-p15-066-enclosing}{The enclosing classes of the earlier capacity failures}
\begin{enumerate}
\item \emph{Searching many features}
      (\cref{thm:p15-conc-selection-example}). With \(K\) predictors and
      \(n\) records, fix \(0<\delta<1\). If
      \(K\geq2^n\ln(1/\delta)\), then with probability at least
      \(1-\delta\) some predictor has empirical risk zero, and on that
      event every exact empirical-risk minimizer \(\widehat h\) satisfies
      \(R(\widehat h)-\widehat R_n(\widehat h)=1/2\).
      If \(K\geq2^n\), the radius \(b_n\) of
      \cref{thm:p15-conc-finite-learning} exceeds \(1/2\) for every
      \(\delta\in(0,1)\), so that theorem remains valid and reports no
      useful bound. If \(\delta\) is fixed and \((\ln K)/n\to0\), then
      \(b_n\to0\).
\item \emph{All binary targets}
      (\cref{thm:p15-occam-unseen-labels}). On an \(N\)-point domain the
      class of all binary functions has dimension \(N\), so
      \cref{thm:p15-w02b-p15-065-lower-realizable} applies to it as
      well.
\item \emph{A class chosen after the data}
      (\cref{thm:p15-occam-boundaries}, item~3). On an \(N\)-point domain
      \(\mathcal X\) with \(N\geq m\) and the uniform input law, the
      smallest fixed class containing the target \(f=0\) and every
      possible output \(h_S\) is
      \[
       H_m=\{f\}\cup\{\mathbf1_{\mathcal X\setminus O}:O\subseteq\mathcal X,\ 1\leq|O|\leq m\}.
      \]
      Its size \(r\) satisfies \(\ln r\geq m\ln(N/m)\), and if
      \(N\geq m+2\) its dimension is \(m+1\). If \(N\geq2m\), the sufficient condition
      of \cref{thm:p15-occam-finite} for this class at sample size \(m\)
      holds only when \(\varepsilon>1/2\), in agreement with the actual
      error at least \(1-m/N\geq1/2\).
\item \emph{Memorizing the sample}
      (\cref{thm:p15-vcl-examples}, item~3). Every classifier the
      memorizer can return has the form
      \(\mathbf1_{[0,1]\setminus F}\) with \(F\subseteq[0,1/2)\) finite.
      The class of all such functions shatters every finite subset of
      \([0,1/2)\), so it has infinite dimension, and
      \cref{thm:p15-w02b-p15-065-lower-infinite} applies to it.
\item \emph{Finite subsets of the integers}
      (\cref{thm:p15-vcl-examples}, item~2). Their indicators form a
      class of infinite dimension; besides the discrepancy failure
      proved there, \cref{thm:p15-w02b-p15-065-lower-infinite}
      gives the corresponding failure of every learner.
\end{enumerate}
\end{hbtheorem}

Items~1 and~3 show how the finite-class theorems behave when applied
honestly to the enclosing class: they remain true and become
uninformative exactly when the actual failure occurs. The failure in
item~1 depends on \(\ln K\) being comparable to \(n\); in item~3 on the
data-dependent class being treated as though it had two members rather
than up to \(1+\sum_{j=1}^m\binom Nj\). Items~4 and~5 concern laws fixed in
advance; the lower bounds of \cref{sec:p15-w02b-p15-065-lower}
instead choose a bad law for each learner, so a failure for some law does not mean failure for
every law, as the point-mass case of \cref{thm:p15-vcl-examples},
item~2, shows.

\section{Distribution shift: transporting a guarantee to a changed law}
\label{sec:p15-w02b-p15-066-shift}

Every guarantee of
\cref{sec:p15-bounded-concentration-learning,sec:p15-occam-learning,sec:p15-vc-dimension-and-uniform-learning}
concerns \(R_P\). To speak about \(R_Q\) one needs an inequality between
the two risks that holds for every candidate, because the selected
candidate is random. The next two theorems supply such inequalities
under two different assumptions about \(Q\).

\begin{hbtheorem}{thm:p15-w02b-p15-066-tv-transport}{proved=annex:p15-w02b-p15-066-tv-transport}{Transport by total variation}
Let \(d=\|P-Q\|_{\mathrm{TV}}\).
\begin{enumerate}
\item For every measurable \(f:\mathcal Z\to[0,1]\),
      \[
       \left|\int f\,dQ-\int f\,dP\right|\leq d .
      \]
      There is \(A\in\mathcal B\) with \(Q(A)-P(A)=d\), so equality holds
      for \(f=\mathbf1_A\). In particular \(|R_Q(h)-R_P(h)|\leq d\) for
      every candidate, with equality for a loss equal to
      \(\mathbf1_A\).
\item In the finite-class setting of \cref{def:p15-conc-learning}, with
      the sample from \(P\), on the event of probability at least
      \(1-\delta\) given by \cref{thm:p15-conc-finite-learning}, every
      \(h\) satisfies \(|\widehat R_n(h)-R_Q(h)|\leq b_n+d\), and every
      \(\eta\)-approximate empirical minimizer satisfies
      \[
       R_Q(\widehat h)\leq\min_{h\in\mathcal H}R_Q(h)+2b_n+2d+\eta .
      \]
      In the setting of \cref{thm:p15-vcl-erm}, on any event where the
      uniform discrepancy under \(P\) is at most \(u\), every
      \(\eta\)-approximate empirical minimizer satisfies
      \(R_Q(\widehat h)\leq\inf_{h\in H}R_Q(h)+2u+2d+\eta\).
      In the realizable setting of \cref{thm:p15-occam-finite}, with
      input laws \(P\) and \(Q\) on \(\mathcal X\) and \(d\) their
      distance, every consistent hypothesis has
      \(R_Q(h,f)\leq\varepsilon+d\) on the event of that theorem.
\item Under covariate shift with a finite label set,
      \(\|P-Q\|_{\mathrm{TV}}=\|P_{\mathcal X}-Q_{\mathcal X}\|_{\mathrm{TV}}\).
\end{enumerate}
\end{hbtheorem}

The term \(d\) is not random and does not shrink with \(n\): more data
from \(P\) reduce \(b_n\) and leave \(d\) untouched. The factor \(2\) in
front of \(d\) in item~2 arises because the selected candidate can lose
\(d\) and the best candidate under \(Q\) can gain \(d\);
\cref{thm:p15-w02b-p15-066-shift-attained} shows both losses occurring
together.

Total variation charges a shift at its largest effect on any event. A
multiplicative assumption is more useful when the risk under \(P\) is
already small, as in the realizable theorems.

\begin{hbtheorem}{thm:p15-w02b-p15-066-density-transport}{proved=annex:p15-w02b-p15-066-density-transport}{Transport by a bounded density ratio}
Suppose there is a measurable \(w:\mathcal Z\to[0,B]\), \(B<\infty\),
with \(Q(A)=\int_Aw\,dP\) for every \(A\in\mathcal B\). Then \(B\geq1\),
and for every measurable \(f\geq0\),
\[
 \int f\,dQ=\int fw\,dP\leq B\int f\,dP .
\]
Hence \(R_Q(h)\leq B\,R_P(h)\) for every candidate. In the setting of
\cref{thm:p15-occam-finite}, with \(Q\) a law on inputs having such a
density relative to \(P\), every consistent hypothesis has
\(R_Q(h,f)\leq B\varepsilon\) on the event of that theorem.
The factor \(B\) cannot be lowered: for \(0<\varepsilon\leq1/B\), the
input laws on \(\{0,1\}\) with \(P\{1\}=\varepsilon\) and
\(Q\{1\}=B\varepsilon\), target \(f=0\) and hypothesis \(h(x)=x\) give
\(R_P(h,f)=\varepsilon\), \(R_Q(h,f)=B\varepsilon\) and a density
bounded by \(B\). If \(Q(A)>0=P(A)\) for some \(A\), no such \(w\)
exists.
\end{hbtheorem}

The next theorem shows that both transfer inequalities are attained, that
the factor \(2\) in the excess-risk transfer is needed, and that without
an assumption on \(Q\) no use of the training sample can certify the
risk under \(Q\).

\begin{hbtheorem}{thm:p15-w02b-p15-066-shift-attained}{proved=annex:p15-w02b-p15-066-shift-attained}{Shift examples that attain the transfer bounds}
\begin{enumerate}
\item \emph{Covariate shift with a consistent hypothesis.} On
      \(\mathcal X=\{0,1\}\) with target \(f=0\) and \(h(x)=x\), let
      \(P=\delta_0\) and \(Q\{1\}=d\), \(0\leq d\leq1\). Then
      \(\|P-Q\|_{\mathrm{TV}}=d\), \(h\) is consistent with every
      training sample, \(R_P(h,f)=0\) and \(R_Q(h,f)=d\). The case
      \(d=1\) is \cref{thm:p15-occam-boundaries}, item~2.
\item \emph{The factor \(2\) for the selected candidate.} Let a single
      input carry a label \(Y\in\{0,1\}\), with \(P\{Y=1\}=1/2+\gamma\)
      and \(Q\{Y=1\}=1/2+\gamma-d\), where \(0<\gamma\leq1/2\) and
      \(\gamma<d\leq1/2+\gamma\). For the constant predictors \(h_1=1\) and
      \(h_2=0\), in that order, \(\|P-Q\|_{\mathrm{TV}}=d\),
      \(h_1\) is the unique minimizer of \(R_P\), and
      \[
       R_Q(h_1)-\min\{R_Q(h_1),R_Q(h_2)\}=2(d-\gamma).
      \]
      The empirical minimizer with this tie order returns \(h_1\) with
      probability at least \(1-e^{-2n\gamma^2}\).
\item \emph{No certificate from the training sample alone.} Let \(h\)
      be a candidate and \(z_1\in\mathcal Z\) a record with
      \(\ell(h,z_1)=1\). For \(0\leq d\leq1\), the law
      \(Q_d=(1-d)P+d\delta_{z_1}\) satisfies
      \(\|P-Q_d\|_{\mathrm{TV}}\leq d\) and
      \(R_{Q_d}(h)=(1-d)R_P(h)+d\). Hence if a function \(T\) of the
      training sample satisfies \(P^{\otimes n}\{R_Q(h)\leq T\}\geq1-\delta\)
      for every \(Q\) with \(\|P-Q\|_{\mathrm{TV}}\leq d\), then
      \(P^{\otimes n}\{T\geq(1-d)R_P(h)+d\}\geq1-\delta\). If the
      requirement is imposed for every law \(Q\), then
      \(P^{\otimes n}\{T\geq1\}\geq1-\delta\).
\item \emph{Total variation ignores geometry.} On the real line let
      \(P=\delta_0\), \(Q=\delta_t\), \(t>0\), with label \(0\), and
      \(h_s(x)=\mathbf1_{\{x\geq s\}}\). Then
      \(\|P-Q\|_{\mathrm{TV}}=1\); \(R_Q(h_s)-R_P(h_s)=1\) for
      \(0<s\leq t\) and \(=0\) otherwise.
\end{enumerate}
\end{hbtheorem}

\begin{hbtheorem}{thm:p15-w02b-p15-066-worked}{proved=annex:p15-w02b-p15-066-worked}{A worked comparison of capacity, covariate shift and concept change}
\begin{enumerate}
\item In \cref{thm:p15-conc-selection-example} take \(n=10\) and
      \(K=3068\geq2^{10}\ln20\). With probability at least \(19/20\)
      some predictor has empirical risk zero, so the empirical minimizer
      has training error \(0\) and risk \(1/2\). For \(\delta=1/20\) the
      finite-class radius is \(b_{10}=\sqrt{\ln(122720)/20}>1/2\).
\item In \cref{thm:p15-conc-noisy-example} the training law has
      \(X\) uniform on \(\{0,1\}\) with
      \(P(Y=1\mid X=0)=1/4\) and \(P(Y=1\mid X=1)=3/4\); with
      \(n=1016\) and \(\delta=1/20\), on the event of
      \cref{thm:p15-conc-finite-learning} the empirical minimizer is the
      predictor \(x\), whose risk under \(P\) is \(1/4\).
      \begin{enumerate}
      \item Under covariate shift to \(Q\{X=1\}=q\), \(0\leq q\leq1\),
            with the same conditional label law, the risks of
            \(0,1,x,1-x\) are \(1/4+q/2,\ 3/4-q/2,\ 1/4,\ 3/4\), and
            \(\|P-Q\|_{\mathrm{TV}}=|q-1/2|\). The selected \(x\) stays
            optimal under every such \(Q\), its risk unchanged, while
            the transfer bound only gives
            \(R_Q(x)\leq1/4+1/10+2|q-1/2|\).
      \item Under the concept change \(Q'(Y=1\mid X=0)=3/4\),
            \(Q'(Y=1\mid X=1)=1/4\), with \(X\) still uniform, the risks
            are \(1/2,1/2,3/4,1/4\) and \(\|P-Q'\|_{\mathrm{TV}}=1/2\).
            The selected \(x\) has \(R_{Q'}(x)-R_P(x)=1/2\), attaining
            the bound of \cref{thm:p15-w02b-p15-066-tv-transport},
            item~1, and excess risk \(1/2\) under \(Q'\).
      \end{enumerate}
\end{enumerate}
\end{hbtheorem}

\section{How the failures relate}
\label{sec:p15-w02b-p15-066-relations}

The capacity failures and the shift failures attack different terms of
one inequality. For a selected \(\widehat h\) and any candidate \(h\),
\[
 R_Q(\widehat h)\leq
 \underbrace{\widehat R_n(\widehat h)}_{\text{fit}}
 +\underbrace{\sup_{g}|R_P(g)-\widehat R_n(g)|}_{\text{capacity}}
 +\underbrace{\|P-Q\|_{\mathrm{TV}}}_{\text{shift}},
\]
where the supremum runs over an enclosing class; this is
\cref{thm:p15-w02b-p15-066-tv-transport}, item~1, combined with the
trivial bound of the selected deviation by the supremum. The capacity
term is controlled by sampling and is the subject of
\cref{thm:p15-conc-finite-learning,thm:p15-vcl-uniform}; its failures
are those of \cref{thm:p15-w02b-p15-066-enclosing}. The shift term is
not a sampling quantity at all: by
\cref{thm:p15-w02b-p15-066-shift-attained}, item~3, it must be bounded
by an assumption, or by data drawn from \(Q\) itself. Independent
validation data from \(Q\) can be used through
\cref{thm:p15-conc-validation} with \(Q\) in place of \(P\), because
that theorem only requires independence from the training object and
the stated law for the validation records. A validation sample from
\(P\) removes the selection effect, because it bounds
\(|R_P(\widehat h)-\widehat R_{\mathrm{val}}(\widehat h)|\) for the
validation average \(\widehat R_{\mathrm{val}}\), but, as the remark after
\cref{thm:p15-vcl-examples} notes, it concerns the law it was drawn
from and leaves the shift term as it was.

High capacity also interacts with shift. In
\cref{thm:p15-w02b-p15-066-shift-attained}, item~4, a threshold between
the supports of \(P\) and \(Q\) turns a tiny shift into a total loss,
while a threshold outside that interval loses nothing. Richer classes
contain more functions that separate the regions where \(P\) and \(Q\)
put their mass, so the bound \(|R_Q(h)-R_P(h)|\leq d\) is attained by
more of their members.

\begin{hbmethod}{def:p15-w02b-p15-066-diagnosis}{Locating a generalization failure}
\begin{enumerate}
\item Write down the selection rule and an enclosing class fixed before
      sampling. If the class was narrowed after the data were seen, use
      the class of all outputs the rule could have produced, as in
      \cref{thm:p15-w02b-p15-066-enclosing}, item~3.
\item Compute the capacity of that class: \(\ln r\) for a finite class,
      the weights of \cref{thm:p15-conc-weighted-learning} for a
      countable one, or the dimension for an infinite one. Compare it
      with the sample size through the applicable theorem. If the
      class shatters many more points than the sample contains, the
      lower bounds of \cref{sec:p15-w02b-p15-065-lower} show that no
      learner can do better without assumptions on the law.
\item Name the evaluation law. If it is the training law, the
      sampling guarantee applies as proved. If not, state a linking
      assumption, a bound on \(\|P-Q\|_{\mathrm{TV}}\) or on a density
      ratio, and add the corresponding term from
      \cref{thm:p15-w02b-p15-066-tv-transport} or
      \cref{thm:p15-w02b-p15-066-density-transport}.
\item If no linking assumption can be justified, obtain an independent
      sample from the evaluation law and use the validation bound of
      \cref{thm:p15-conc-validation}; otherwise report that no
      guarantee for the evaluation law has been established.
\end{enumerate}
\end{hbmethod}

\begin{hbwarning}{def:p15-w02b-p15-066-limits}{What these results do not establish}
The capacity diagnosis of \cref{thm:p15-w02b-p15-066-enclosing} shows
that a guarantee is uninformative, not that learning fails for the law at
hand. For a fixed law a high-capacity class may still be learned, as the
point-mass case of \cref{thm:p15-vcl-examples}, item~2, shows; the lower
bounds that hold for every learner choose the law after the learner, as
\cref{def:p15-w02b-p15-065-minimax-warning} explains.

The total-variation transfer is sharp as a statement over all losses with
values in \([0,1]\), but for a particular loss it can be loose by a large
factor, as in the covariate shift of
\cref{thm:p15-w02b-p15-066-worked}, item~2(a). It equals one for laws
with disjoint supports however close those supports are. The density-ratio
transfer requires that \(Q\) give no mass to events that are null under
\(P\). The training sample has the same law whatever \(Q\) is, so it
carries no information about either quantity;
\cref{thm:p15-w02b-p15-066-shift-attained}, item~3, turns this into the
impossibility of a certificate for \(R_Q\) without an assumption on
\(Q\).
Transfer under other distances between laws, and methods that reweight
the training data toward \(Q\), need their own assumptions and theorems
and are not treated here.
\end{hbwarning}
